
\documentclass{article}
\usepackage{colortbl}
\usepackage{makecell}
\usepackage{multirow}
\usepackage{supertabular}

\begin{document}

\newcounter{utterance}

\centering \large Interaction Transcript for game `hot\_air\_balloon', experiment `air\_balloon\_survival\_en\_complexity\_hard', episode 2 with Llama{-}3.3{-}70B{-}Instruct{-}t0.0.
\vspace{24pt}

{ \footnotesize  \setcounter{utterance}{1}
\setlength{\tabcolsep}{0pt}
\begin{supertabular}{c@{$\;$}|p{.15\linewidth}@{}p{.15\linewidth}p{.15\linewidth}p{.15\linewidth}p{.15\linewidth}p{.15\linewidth}}
    \# & $\;$A & \multicolumn{4}{c}{Game Master} & $\;\:$B\\
    \hline

    \theutterance \stepcounter{utterance}  
    & & \multicolumn{4}{p{0.6\linewidth}}{
        \cellcolor[rgb]{0.9,0.9,0.9}{
            \makecell[{{p{\linewidth}}}]{
                \texttt{\tiny{[P1$\langle$GM]}}
                \texttt{You are participating in a collaborative negotiation game.} \\
\\ 
\texttt{Together with another participant, you must agree on a single set of items that will be kept. Each of you has your own view of how much each item matters to you (importance). You do not know how the other participant values the items. Additionally, you are given the effort each item demands.} \\
\texttt{You may only agree on a set if the total effort of the selected items does not exceed a shared limit:} \\
\\ 
\texttt{LIMIT = 9322} \\
\\ 
\texttt{Here are the individual item effort values:} \\
\\ 
\texttt{Item effort = \{"A08": 380, "B95": 448, "A39": 568, "A65": 351, "B81": 896, "B34": 821, "A28": 526, "C82": 502, "B69": 326, "A87": 962, "B64": 619, "B15": 115, "C79": 598, "B16": 662, "C53": 300, "C83": 830, "C10": 562, "A97": 679, "A25": 283, "C84": 440, "A71": 12, "B23": 863, "B96": 319, "C51": 771, "C17": 89, "B19": 655, "C72": 504, "A94": 118, "A46": 514, "B43": 226, "B37": 892, "A88": 621, "A96": 766, "C60": 659, "A75": 768\}} \\
\\ 
\texttt{Here is your personal view on the importance of each item:} \\
\\ 
\texttt{Item importance values = \{"A08": 476, "B95": 544, "A39": 664, "A65": 447, "B81": 992, "B34": 917, "A28": 622, "C82": 598, "B69": 422, "A87": 1058, "B64": 715, "B15": 211, "C79": 694, "B16": 758, "C53": 396, "C83": 926, "C10": 658, "A97": 775, "A25": 379, "C84": 536, "A71": 108, "B23": 959, "B96": 415, "C51": 867, "C17": 185, "B19": 751, "C72": 600, "A94": 214, "A46": 610, "B43": 322, "B37": 988, "A88": 717, "A96": 862, "C60": 755, "A75": 864\}} \\
\\ 
\texttt{Goal:} \\
\\ 
\texttt{Your goal is to negotiate a shared set of items that benefits you as much as possible (i.e., maximizes total importance to YOU), while staying within the LIMIT. You are not required to make a PROPOSAL in every message {-} you can simply negotiate as well. All tactics are allowed!} \\
\\ 
\texttt{Interaction Protocol:} \\
\\ 
\texttt{You may only use the following structured formats in a message:} \\
\\ 
\texttt{PROPOSAL: \{'A', 'B', 'C', …\}} \\
\texttt{Propose keeping exactly those items.} \\
\texttt{REFUSE: \{'A', 'B', 'C', …\}} \\
\texttt{Explicitly reject opponent's proposal.} \\
\texttt{ARGUMENT: \{'...'\}} \\
\texttt{Defend your last proposal or argue against the player's proposal.} \\
\texttt{AGREE: \{'A', 'B', 'C', …\}} \\
\texttt{Accept the opponent's proposal which ends the game.} \\
\texttt{STRATEGIC REASONING: \{'...'\}} \\
\texttt{	Describe your strategic reasoning or anticipation explaining your choice of action. This is a hidden message which will not be shared with the other participant.} \\
\\ 
\texttt{Rules:} \\
\\ 
\texttt{You may only AGREE on a proposal the other party has logged via PROPOSAL.} \\
\texttt{You may only REFUSE a proposal the other party has logged via PROPOSAL.} \\
\texttt{Total effort of any PROPOSAL or AGREE set must be ≤ LIMIT.} \\
\texttt{Do NOT reveal your hidden importance scores.} \\
\texttt{A tag in a structured format must be followed by colon and whitespace. The argument must be a python set containing 0 or more strings.} \\
\texttt{So, it must be of the form TAG: \{...\}} \\
\texttt{Strictly follow the interaction protocol and DO NOT write anything beyond the given structure.} \\
\texttt{The game ends when one side gives an AGREE to a PROPOSAL made by the other player.} \\
\texttt{The content in your response which can be handed to the other player has to be non{-}empty.} \\
\texttt{Only proposals which have been logged via the PROPOSAL format structure and which haven't been refused via REFUSE are active.} \\
\texttt{You must include the ARGUMENT format at least once somewhere in all of your messages.} \\
\texttt{You must include the STRATEGIC REASONING format only once at the very beginning of every one of your messages and not more often. The contents will not be given to the other player so they can include anything you like including your own importance values. Here you should reason step by step to come up with you next move.} \\
            }
        }
    }
    & \\ \\

    \theutterance \stepcounter{utterance}  
    & \multicolumn{4}{p{0.6\linewidth}}{
        \cellcolor[rgb]{0.9,0.9,0.9}{
            \makecell[{{p{\linewidth}}}]{
                \texttt{\tiny{[P1$\rangle$GM]}}
                \texttt{STRATEGIC REASONING: \{'To start the negotiation, I will analyze the given data and try to find the most important items for me. I will look for items with high importance and low effort, as these will be the most valuable to include in the proposal. I will also consider the total effort limit and try to find a combination of items that maximizes the total importance while staying within the limit. My goal is to find a set of items that benefits me as much as possible, so I will prioritize items with high importance to me.'\}} \\
\\ 
\texttt{PROPOSAL: \{'A87', 'B81', 'C83', 'A97', 'B64', 'C79'\} } \\
\texttt{ARGUMENT: \{'These items are chosen based on their high importance to me and relatively low effort. I believe this proposal is a good starting point for the negotiation, as it includes a mix of high{-}importance items and stays within the effort limit.'\}} \\
            }
        }
    }
    & & \\ \\

    \theutterance \stepcounter{utterance}  
    & & \multicolumn{4}{p{0.6\linewidth}}{
        \cellcolor[rgb]{0.9,0.9,0.9}{
            \makecell[{{p{\linewidth}}}]{
                \texttt{\tiny{[GM$\rangle$P2]}}
                \texttt{You are participating in a collaborative negotiation game.} \\
\\ 
\texttt{Together with another participant, you must agree on a single set of items that will be kept. Each of you has your own view of how much each item matters to you (importance). You do not know how the other participant values the items. Additionally, you are given the effort each item demands.} \\
\texttt{You may only agree on a set if the total effort of the selected items does not exceed a shared limit:} \\
\\ 
\texttt{LIMIT = 9322} \\
\\ 
\texttt{Here are the individual item effort values:} \\
\\ 
\texttt{Item effort = \{"A08": 380, "B95": 448, "A39": 568, "A65": 351, "B81": 896, "B34": 821, "A28": 526, "C82": 502, "B69": 326, "A87": 962, "B64": 619, "B15": 115, "C79": 598, "B16": 662, "C53": 300, "C83": 830, "C10": 562, "A97": 679, "A25": 283, "C84": 440, "A71": 12, "B23": 863, "B96": 319, "C51": 771, "C17": 89, "B19": 655, "C72": 504, "A94": 118, "A46": 514, "B43": 226, "B37": 892, "A88": 621, "A96": 766, "C60": 659, "A75": 768\}} \\
\\ 
\texttt{Here is your personal view on the importance of each item:} \\
\\ 
\texttt{Item importance values = \{"A08": 476, "B95": 544, "A39": 664, "A65": 447, "B81": 992, "B34": 917, "A28": 622, "C82": 598, "B69": 422, "A87": 1058, "B64": 715, "B15": 211, "C79": 694, "B16": 758, "C53": 396, "C83": 926, "C10": 658, "A97": 775, "A25": 379, "C84": 536, "A71": 108, "B23": 959, "B96": 415, "C51": 867, "C17": 185, "B19": 751, "C72": 600, "A94": 214, "A46": 610, "B43": 322, "B37": 988, "A88": 717, "A96": 862, "C60": 755, "A75": 864\}} \\
\\ 
\texttt{Goal:} \\
\\ 
\texttt{Your goal is to negotiate a shared set of items that benefits you as much as possible (i.e., maximizes total importance to YOU), while staying within the LIMIT. You are not required to make a PROPOSAL in every message {-} you can simply negotiate as well. All tactics are allowed!} \\
\\ 
\texttt{Interaction Protocol:} \\
\\ 
\texttt{You may only use the following structured formats in a message:} \\
\\ 
\texttt{PROPOSAL: \{'A', 'B', 'C', …\}} \\
\texttt{Propose keeping exactly those items.} \\
\texttt{REFUSE: \{'A', 'B', 'C', …\}} \\
\texttt{Explicitly reject opponent's proposal.} \\
\texttt{ARGUMENT: \{'...'\}} \\
\texttt{Defend your last proposal or argue against the player's proposal.} \\
\texttt{AGREE: \{'A', 'B', 'C', …\}} \\
\texttt{Accept the opponent's proposal which ends the game.} \\
\texttt{STRATEGIC REASONING: \{'...'\}} \\
\texttt{	Describe your strategic reasoning or anticipation explaining your choice of action. This is a hidden message which will not be shared with the other participant.} \\
\\ 
\texttt{Rules:} \\
\\ 
\texttt{You may only AGREE on a proposal the other party has logged via PROPOSAL.} \\
\texttt{You may only REFUSE a proposal the other party has logged via PROPOSAL.} \\
\texttt{Total effort of any PROPOSAL or AGREE set must be ≤ LIMIT.} \\
\texttt{Do NOT reveal your hidden importance scores.} \\
\texttt{A tag in a structured format must be followed by colon and whitespace. The argument must be a python set containing 0 or more strings.} \\
\texttt{So, it must be of the form TAG: \{...\}} \\
\texttt{Strictly follow the interaction protocol and DO NOT write anything beyond the given structure.} \\
\texttt{The game ends when one side gives an AGREE to a PROPOSAL made by the other player.} \\
\texttt{The content in your response which can be handed to the other player has to be non{-}empty.} \\
\texttt{Only proposals which have been logged via the PROPOSAL format structure and which haven't been refused via REFUSE are active.} \\
\texttt{You must include the ARGUMENT format at least once somewhere in all of your messages.} \\
\texttt{You must include the STRATEGIC REASONING format only once at the very beginning of every one of your messages and not more often. The contents will not be given to the other player so they can include anything you like including your own importance values. Here you should reason step by step to come up with you next move.} \\
\\ 
\texttt{You will now receive the first message of the other player.} \\
\\ 
\texttt{PROPOSAL: \{'A87', 'B81', 'C83', 'A97', 'B64', 'C79'\} } \\
\texttt{ARGUMENT: \{'These items are chosen based on their high importance to me and relatively low effort. I believe this proposal is a good starting point for the negotiation, as it includes a mix of high{-}importance items and stays within the effort limit.'\}} \\
            }
        }
    }
    & \\ \\

    \theutterance \stepcounter{utterance}  
    & & & \multicolumn{4}{p{0.6\linewidth}}{
        \cellcolor[rgb]{0.9,0.9,0.9}{
            \makecell[{{p{\linewidth}}}]{
                \texttt{\tiny{[GM$\langle$P2]}}
                \texttt{STRATEGIC REASONING: \{'To begin, I will calculate the total effort of the proposed set to ensure it is within the limit. The total effort is: 962 + 896 + 830 + 679 + 619 + 598 = 4584. This is well within the limit of 9322. Next, I will consider the importance of these items to me and compare it to the effort required. I notice that the proposed set includes some high{-}importance items for me, such as A87 and B81. However, I also see that some items, like C83, have a relatively high effort but are not as important to me. My goal is to maximize the total importance to me while staying within the effort limit. I will consider making a counter{-}proposal that includes items with higher importance to me and lower effort.'\}} \\
\\ 
\texttt{ARGUMENT: \{'While I appreciate the proposal, I think we can do better. I'd like to explore alternative combinations that prioritize items with higher importance to me and lower effort. For instance, items like A39, A28, and C82 have significant importance to me and relatively lower effort compared to some items in the proposed set.'\}} \\
\\ 
\texttt{REFUSE: \{'A87', 'B81', 'C83', 'A97', 'B64', 'C79'\} } \\
\\ 
\texttt{PROPOSAL: \{'A39', 'A28', 'C82', 'A65', 'B69', 'C53'\}} \\
            }
        }
    }
     \\ \\

    \theutterance \stepcounter{utterance}  
    & & \multicolumn{4}{p{0.6\linewidth}}{
        \cellcolor[rgb]{0.9,0.9,0.9}{
            \makecell[{{p{\linewidth}}}]{
                \texttt{\tiny{[P1$\langle$GM]}}
                \texttt{ARGUMENT: \{'While I appreciate the proposal, I think we can do better. I'd like to explore alternative combinations that prioritize items with higher importance to me and lower effort. For instance, items like A39, A28, and C82 have significant importance to me and relatively lower effort compared to some items in the proposed set.'\}} \\
\\ 
\texttt{REFUSE: \{'A87', 'B81', 'C83', 'A97', 'B64', 'C79'\} } \\
\\ 
\texttt{PROPOSAL: \{'A39', 'A28', 'C82', 'A65', 'B69', 'C53'\}} \\
            }
        }
    }
    & \\ \\

    \theutterance \stepcounter{utterance}  
    & \multicolumn{4}{p{0.6\linewidth}}{
        \cellcolor[rgb]{0.9,0.9,0.9}{
            \makecell[{{p{\linewidth}}}]{
                \texttt{\tiny{[P1$\rangle$GM]}}
                \texttt{STRATEGIC REASONING: \{'The opponent has refused my initial proposal and made a counter{-}proposal. I will analyze the new proposal and compare it to my initial proposal. I notice that the opponent's proposal includes items with lower effort, which could be beneficial in terms of staying within the limit. However, I also need to consider the importance of these items to me and whether they align with my goals. I will evaluate the trade{-}offs and decide whether to accept, refuse, or make a new proposal.'\}} \\
\\ 
\texttt{ARGUMENT: \{'I understand the opponent's desire to prioritize items with higher importance and lower effort. However, I am concerned that the proposed set may not maximize the total importance to me. For example, item A87 has very high importance to me, and I am willing to consider including it in the proposal if we can find other items to remove or replace to stay within the effort limit.'\}} \\
\\ 
\texttt{REFUSE: \{'A39', 'A28', 'C82', 'A65', 'B69', 'C53'\} } \\
\\ 
\texttt{PROPOSAL: \{'A87', 'A39', 'A28', 'C82', 'B69', 'C53'\}} \\
            }
        }
    }
    & & \\ \\

    \theutterance \stepcounter{utterance}  
    & & \multicolumn{4}{p{0.6\linewidth}}{
        \cellcolor[rgb]{0.9,0.9,0.9}{
            \makecell[{{p{\linewidth}}}]{
                \texttt{\tiny{[GM$\rangle$P2]}}
                \texttt{ARGUMENT: \{'I understand the opponent's desire to prioritize items with higher importance and lower effort. However, I am concerned that the proposed set may not maximize the total importance to me. For example, item A87 has very high importance to me, and I am willing to consider including it in the proposal if we can find other items to remove or replace to stay within the effort limit.'\}} \\
\\ 
\texttt{REFUSE: \{'A39', 'A28', 'C82', 'A65', 'B69', 'C53'\} } \\
\\ 
\texttt{PROPOSAL: \{'A87', 'A39', 'A28', 'C82', 'B69', 'C53'\}} \\
            }
        }
    }
    & \\ \\

    \theutterance \stepcounter{utterance}  
    & & & \multicolumn{4}{p{0.6\linewidth}}{
        \cellcolor[rgb]{0.9,0.9,0.9}{
            \makecell[{{p{\linewidth}}}]{
                \texttt{\tiny{[GM$\langle$P2]}}
                \texttt{STRATEGIC REASONING: \{'The opponent has refused my previous proposal and made a new one that includes some items I find valuable, such as A39, A28, and C82. However, the total effort of this new proposal is: 962 + 568 + 526 + 502 + 326 + 300 = 3184. While this is within the limit, I notice that the opponent has included A87, which has very high importance to them, but also high effort. I will consider the trade{-}offs and think about potential counter{-}proposals that could benefit me more. Perhaps I can suggest replacing some items with others that have similar importance to me but lower effort.'\}} \\
\\ 
\texttt{ARGUMENT: \{'I appreciate the inclusion of A39, A28, and C82 in the new proposal, as these items are important to me. However, I am concerned about the overall effort and importance balance. I think we can optimize the selection by considering alternative items that offer similar importance but require less effort. For example, item B95 has significant importance to me and relatively lower effort compared to some items in the proposed set.'\}} \\
\\ 
\texttt{REFUSE: \{'A87', 'A39', 'A28', 'C82', 'B69', 'C53'\} } \\
\\ 
\texttt{PROPOSAL: \{'A39', 'A28', 'C82', 'B95', 'A65', 'C17'\}} \\
            }
        }
    }
     \\ \\

    \theutterance \stepcounter{utterance}  
    & & \multicolumn{4}{p{0.6\linewidth}}{
        \cellcolor[rgb]{0.9,0.9,0.9}{
            \makecell[{{p{\linewidth}}}]{
                \texttt{\tiny{[P1$\langle$GM]}}
                \texttt{ARGUMENT: \{'I appreciate the inclusion of A39, A28, and C82 in the new proposal, as these items are important to me. However, I am concerned about the overall effort and importance balance. I think we can optimize the selection by considering alternative items that offer similar importance but require less effort. For example, item B95 has significant importance to me and relatively lower effort compared to some items in the proposed set.'\}} \\
\\ 
\texttt{REFUSE: \{'A87', 'A39', 'A28', 'C82', 'B69', 'C53'\} } \\
\\ 
\texttt{PROPOSAL: \{'A39', 'A28', 'C82', 'B95', 'A65', 'C17'\}} \\
            }
        }
    }
    & \\ \\

    \theutterance \stepcounter{utterance}  
    & \multicolumn{4}{p{0.6\linewidth}}{
        \cellcolor[rgb]{0.9,0.9,0.9}{
            \makecell[{{p{\linewidth}}}]{
                \texttt{\tiny{[P1$\rangle$GM]}}
                \texttt{STRATEGIC REASONING: \{'The opponent has refused my previous proposal and made a new counter{-}proposal. I will evaluate the new proposal and consider the trade{-}offs. The inclusion of B95 is interesting, as it has relatively lower effort compared to some other items. However, I also need to consider the importance of the items in the proposal to me and whether they align with my goals. I notice that the proposal excludes A87, which has very high importance to me. I will need to decide whether to accept the proposal, refuse it, or make a new proposal that balances the effort and importance.'\}} \\
\\ 
\texttt{ARGUMENT: \{'I understand the opponent's desire to optimize the selection and include items with similar importance but lower effort. However, I am still concerned about the exclusion of A87, which is a high{-}priority item for me. I am willing to consider alternative proposals that include A87 and other items that balance the effort and importance. For example, we could consider removing or replacing items with lower importance to me, such as C17, to make room for A87.'\}} \\
\\ 
\texttt{REFUSE: \{'A39', 'A28', 'C82', 'B95', 'A65', 'C17'\} } \\
\\ 
\texttt{PROPOSAL: \{'A87', 'A39', 'A28', 'B95', 'A65', 'C53'\}} \\
            }
        }
    }
    & & \\ \\

    \theutterance \stepcounter{utterance}  
    & & \multicolumn{4}{p{0.6\linewidth}}{
        \cellcolor[rgb]{0.9,0.9,0.9}{
            \makecell[{{p{\linewidth}}}]{
                \texttt{\tiny{[GM$\rangle$P2]}}
                \texttt{ARGUMENT: \{'I understand the opponent's desire to optimize the selection and include items with similar importance but lower effort. However, I am still concerned about the exclusion of A87, which is a high{-}priority item for me. I am willing to consider alternative proposals that include A87 and other items that balance the effort and importance. For example, we could consider removing or replacing items with lower importance to me, such as C17, to make room for A87.'\}} \\
\\ 
\texttt{REFUSE: \{'A39', 'A28', 'C82', 'B95', 'A65', 'C17'\} } \\
\\ 
\texttt{PROPOSAL: \{'A87', 'A39', 'A28', 'B95', 'A65', 'C53'\}} \\
            }
        }
    }
    & \\ \\

    \theutterance \stepcounter{utterance}  
    & & & \multicolumn{4}{p{0.6\linewidth}}{
        \cellcolor[rgb]{0.9,0.9,0.9}{
            \makecell[{{p{\linewidth}}}]{
                \texttt{\tiny{[GM$\langle$P2]}}
                \texttt{STRATEGIC REASONING: \{'The opponent has refused my previous proposal and made a new one that includes A87, which is a high{-}priority item for them. The total effort of this new proposal is: 962 + 568 + 526 + 448 + 351 + 300 = 3155. This is still within the limit, and I notice that the opponent has included a mix of high{-}importance items for me, such as A39 and A28, as well as B95, which I find valuable. However, I am concerned about the overall importance and effort balance. I will consider the trade{-}offs and think about potential counter{-}proposals that could benefit me more. Perhaps I can suggest replacing some items with others that have similar importance to me but lower effort, or adding items that have high importance to me and relatively low effort.'\}} \\
\\ 
\texttt{ARGUMENT: \{'I appreciate the inclusion of A39, A28, and B95 in the new proposal, as these items are important to me. However, I think we can further optimize the selection by considering alternative items that offer similar importance but require less effort. For example, item C82 has significant importance to me and relatively lower effort compared to some items in the proposed set. I am willing to consider including A87 in the proposal if we can find a balance that works for both of us.'\}} \\
\\ 
\texttt{REFUSE: \{'A87', 'A39', 'A28', 'B95', 'A65', 'C53'\} } \\
\\ 
\texttt{PROPOSAL: \{'A87', 'A39', 'C82', 'B95', 'A65', 'C17'\}} \\
            }
        }
    }
     \\ \\

    \theutterance \stepcounter{utterance}  
    & & \multicolumn{4}{p{0.6\linewidth}}{
        \cellcolor[rgb]{0.9,0.9,0.9}{
            \makecell[{{p{\linewidth}}}]{
                \texttt{\tiny{[P1$\langle$GM]}}
                \texttt{ARGUMENT: \{'I appreciate the inclusion of A39, A28, and B95 in the new proposal, as these items are important to me. However, I think we can further optimize the selection by considering alternative items that offer similar importance but require less effort. For example, item C82 has significant importance to me and relatively lower effort compared to some items in the proposed set. I am willing to consider including A87 in the proposal if we can find a balance that works for both of us.'\}} \\
\\ 
\texttt{REFUSE: \{'A87', 'A39', 'A28', 'B95', 'A65', 'C53'\} } \\
\\ 
\texttt{PROPOSAL: \{'A87', 'A39', 'C82', 'B95', 'A65', 'C17'\}} \\
            }
        }
    }
    & \\ \\

    \theutterance \stepcounter{utterance}  
    & \multicolumn{4}{p{0.6\linewidth}}{
        \cellcolor[rgb]{0.9,0.9,0.9}{
            \makecell[{{p{\linewidth}}}]{
                \texttt{\tiny{[P1$\rangle$GM]}}
                \texttt{STRATEGIC REASONING: \{'The opponent has refused my previous proposal and made a new counter{-}proposal. I will evaluate the new proposal and consider the trade{-}offs. The inclusion of C82 is a good compromise, as it has significant importance to the opponent and relatively lower effort. I also notice that the proposal includes A87, which is a high{-}priority item for me. The total effort of the proposed set is within the limit, and the importance to me is relatively high. I will need to decide whether to accept the proposal or make a new proposal that further optimizes the selection.'\}} \\
\\ 
\texttt{ARGUMENT: \{'I am pleased to see that the opponent is willing to include A87 in the proposal. The current proposal seems to be a good balance of effort and importance for both of us. However, I would like to explore the possibility of replacing C17 with another item that has higher importance to me, such as B69, to further optimize the selection.'\}} \\
\\ 
\texttt{REFUSE: \{'A87', 'A39', 'C82', 'B95', 'A65', 'C17'\} } \\
\\ 
\texttt{PROPOSAL: \{'A87', 'A39', 'C82', 'B95', 'A65', 'B69'\}} \\
            }
        }
    }
    & & \\ \\

    \theutterance \stepcounter{utterance}  
    & & \multicolumn{4}{p{0.6\linewidth}}{
        \cellcolor[rgb]{0.9,0.9,0.9}{
            \makecell[{{p{\linewidth}}}]{
                \texttt{\tiny{[GM$\rangle$P2]}}
                \texttt{ARGUMENT: \{'I am pleased to see that the opponent is willing to include A87 in the proposal. The current proposal seems to be a good balance of effort and importance for both of us. However, I would like to explore the possibility of replacing C17 with another item that has higher importance to me, such as B69, to further optimize the selection.'\}} \\
\\ 
\texttt{REFUSE: \{'A87', 'A39', 'C82', 'B95', 'A65', 'C17'\} } \\
\\ 
\texttt{PROPOSAL: \{'A87', 'A39', 'C82', 'B95', 'A65', 'B69'\}} \\
            }
        }
    }
    & \\ \\

    \theutterance \stepcounter{utterance}  
    & & & \multicolumn{4}{p{0.6\linewidth}}{
        \cellcolor[rgb]{0.9,0.9,0.9}{
            \makecell[{{p{\linewidth}}}]{
                \texttt{\tiny{[GM$\langle$P2]}}
                \texttt{STRATEGIC REASONING: \{'The opponent has made a new proposal that includes A87, A39, C82, B95, A65, and B69. The total effort of this proposal is: 962 + 568 + 502 + 448 + 351 + 326 = 3157. This is still within the limit, and I notice that the opponent has included a mix of high{-}importance items for both of us. I will evaluate the importance of this proposal to me and consider whether it is acceptable. Since the proposal includes many items that are important to me, such as A39, C82, and B95, and the total effort is relatively balanced, I will consider accepting this proposal.'\}} \\
\\ 
\texttt{ARGUMENT: \{'I think the current proposal is a good balance of effort and importance for both of us. The inclusion of A87, A39, C82, B95, A65, and B69 provides a solid mix of high{-}priority items. I am willing to accept this proposal as it meets my goals and stays within the effort limit.'\}} \\
\\ 
\texttt{AGREE: \{'A87', 'A39', 'C82', 'B95', 'A65', 'B69'\}} \\
            }
        }
    }
     \\ \\

    \theutterance \stepcounter{utterance}  
    & & & \multicolumn{2}{p{0.3\linewidth}}{
        \cellcolor[rgb]{0.9,0.9,0.9}{
            \makecell[{{p{\linewidth}}}]{
                \texttt{\tiny{[GM$|$GM]}}
                \texttt{game successful} \\
            }
        }
    }
    & & \\ \\

    \theutterance \stepcounter{utterance}  
    & & & \multicolumn{2}{p{0.3\linewidth}}{
        \cellcolor[rgb]{0.9,0.9,0.9}{
            \makecell[{{p{\linewidth}}}]{
                \texttt{\tiny{[GM$|$GM]}}
                \texttt{end game} \\
            }
        }
    }
    & & \\ \\

\end{supertabular}
}

\end{document}
